% Generated by report/tools/tables.py from the result files of the repository. Do not edit.
\begingroup
\footnotesize
\setlength{\tabcolsep}{3.5pt}
\setlength{\LTcapwidth}{\textwidth}
\begin{longtable}{@{}r>{\raggedright\arraybackslash}p{3.1cm}>{\raggedright\arraybackslash}p{10.95cm}@{}}
\caption[Deviations from the three protocols]{The 89 deviations from the three protocols, in the order of their lists: survey (\file{study/PROTOCOL-en.md}), experiments (\file{sim/SIM\_PROTOCOL-en.md}), GRPO strand (\file{sim/llm/LLM\_PROTOCOL-en.md}). Timing: what the protocol declares had already been seen when the deviation was written (before or after the results concerned, or no effect on the results for the organisation of the computation). The summary takes one line; the full text, which prevails, is in the protocol (the French files are the dated originals). The number of deviations and their numbering are checked against the files by \file{report/tools/tables.py}.}\label{tab:ecarts}\\
\toprule
No. & Timing & Summary \\
\midrule
\endfirsthead
\multicolumn{3}{l}{\emph{Table \thetable, continued}}\\
\toprule
No. & Timing & Summary \\
\midrule
\endhead
\midrule\multicolumn{3}{r}{\emph{continued on next page}}\\
\endfoot
\bottomrule
\endlastfoot
\multicolumn{3}{@{}l}{\textbf{Survey} (\file{study/PROTOCOL-en.md}, 35 deviations)} \\
1 & during coding, before aggregation & X4 harmonised: \texttt{yes} requires a described tuning of the baselines; `same configuration' with no tuning becomes \texttt{no}. \\
2 & during coding, before aggregation & X7 harmonised: \texttt{yes} requires a quantity of compute (hours, GPU-hours, FLOP); the hardware name alone is \texttt{no}. \\
3 & during coding, before aggregation & Papers with no experiment (Q1 \texttt{no}): Q2 to Q8 and X1 to X4 set to \texttt{na}, out of the denominators. \\
4 & during coding, before aggregation & Q7 for a theory paper with toy experiments: \texttt{no}, not \texttt{na}. \\
5 & during coding, before aggregation & X3: a selection unstated but visible in the paper's own tables (LR10) counts as \texttt{yes}. \\
6 & during coding, before aggregation & Q5: code `to be made public' is \texttt{no}; a project page is \texttt{yes} only if it leads to the code. \\
7 & during coding, before aggregation & Calibration on 12 papers from 2018: our grid stricter than Pineau's on Q3, more lenient on Q7; a change must hold against both references. \\
8 & during coding, before aggregation & Separate 2018 appendices read (ICML, NeurIPS); in 2026 they are inside the PDF; C21 has no published appendix. \\
9 & during coding, before aggregation & Contact address removed from the User-Agent header of the downloads; no effect on the data. \\
10 & during coding, before aggregation & Working scripts overwritten in the readers' shared space; the codings, in separate folders, are untouched. \\
11 & before the results & Grid v2 (Y1 to Y12) for the flaws grid v1 leaves unmeasured; Y11 and Y12 added next, before any coding. \\
12 & before the results & Enlargement: 2018 brought to 50 papers (A2), classical RL 2026 to 50 (A3, title screen); analysis fixed (Newcombe, Fisher, Holm); double reading of 10 papers. \\
13 & before the results & Blind human review of 10 papers planned (A1), abandoned in deviation 23. \\
14 & before the results & Decision rules of grid v2 (scope, edge cases of Y1 to Y12); sheet \file{GRILLE.md}, check \file{check\_codes.py}. \\
15 & after the results & Rules of deviations 1, 2 and 4 applied to the 12 papers from 2018 as well, left out until then; no new rule. \\
16 & before the results & Pilot of three papers (LR08, ML15, C218) to find ambiguous rules; its readings enter no analysis. \\
17 & before the results & Rules clarified after the pilot (2018 appendices, Y1, Y3, Y5 to Y9, Y11, format); 9 unreachable ICLR 2026 papers skipped in the A3 queue. \\
18 & after reading A3, before any percentage & X5 (family) read on the model the RL trains, to decide which A3 papers enter the classical sample. \\
19 & after the audit of the readings, before any percentage & Harmonisation of grid v2: 80 codes changed, each with its reason in \file{overrides\_v2.json}, for both years. \\
20 & after the results & Verdict rule: only the same-grid comparison (classical RL 2026 against our reading of 2018) decides; Pineau's figures remain context. \\
21 & after the results, before this computation & Exploratory checks (provenance, PDF source, Mantel-Haenszel, reader removed) and blind test of the cost and risk thesis. \\
22 & before the results & Checklists of the 30 NeurIPS 2025 papers read, extracted and set against our codes (A5). \\
23 & before the results & Human review (A1) abandoned, by our own decision; stated among the limitations. \\
24 & before the results & Sensitivity to the title filter (A6): 400 papers outside the filter, abstracts read, then 20 hidden RL papers read with both grids. \\
25 & before the results & Trajectory from 2018 to 2026 (A4): 20 classical RL papers per year from ICML and NeurIPS 2019, 2020, 2022 and 2024; Cochran-Armitage. \\
26 & before the results & Blind second reading of 10 A4 papers (reader R3). \\
27 & after the results, before any search & Keyword search against the `no' decisions of the main comparison, blind arbiter; decisions turned, then doubtful ones. \\
28 & after the results & Provenance of the readers published (\file{readers/}): two readers finished on \texttt{claude-opus-5} (13 papers); the readings do not change. \\
29 & after the results, before this computation & Main comparison and trajectory redone without the 13 papers read partly by \texttt{claude-opus-5}. \\
30 & after the results, before this computation & Further analyses with no new reading: intervals of the agreement, minimum detectable difference, original verdict rule, readers and year confounded. \\
31 & after the results, before any download & Q5 reread on the latest arXiv version (2018) and on the proceedings (2026), blind arbiter, creation dates of the repositories. \\
32 & after the results, before any reading & New criterion X8: uncertainty computed over seeds, read by eight fresh agents across both years. \\
33 & after the results, before any reading & A reader from another model family (\texttt{gemma3:12b}, local) on the 98 papers of the main comparison. \\
35 & after the results, before any draw and any human reading & Blind human re-reading of both years by us: 15 papers from 2018 and 15 from 2026, criteria Q5, Y9, X7, X2, X8. \\
36 & after the results, before any reading and any computation & Year and readers disentangled: eight fresh readers each read both years on Q5, Y9, X7 and X2. \\
\addlinespace
\multicolumn{3}{@{}l}{\textbf{Experiments} (\file{sim/SIM\_PROTOCOL-en.md}, 34 deviations)} \\
1 & before the first training run & SAC and CrossQ: 200,000 steps instead of 300,000, for lack of compute; TD3 not launched. \\
2 & after the results & E6: an error on our side (normalisation gamma 0.99 instead of 0.999); corrected recipe rerun, diagnosis on Hopper v4 and v5. \\
3 & after the results & E1: a second statistic, the mean $t$ along the curve as in Henderson et al. (\file{e1\_curve.py}). \\
4 & before any TRPO run & B1: 30 TRPO runs on HalfCheetah and three readings of Henderson et al.'s statistic. \\
5 & after \file{e6\_diag}, before any official run & E6: ten agents trained with the Zoo's official script, evaluated with our protocol. \\
6 & after \file{b1\_trpo}, before this computation & B1: the script Henderson published in 2017 (\file{from\_same\_dist.py}) applied to our curves. \\
7 & after PPO and SAC, before looking at CrossQ & Analysis details of E2, E4 and E5 (final score, IQM, probability of improvement, reversals, draws). \\
8 & after PPO and SAC, before looking at CrossQ & B3: three measures (final score, best score, mean of the curve) and ranking reversals. \\
9 & no effect on the results & Organisation: the second worker takes the CrossQ queue from the other end. \\
10 & no effect on the results & Stop by the watchdog (swap); two CrossQ runs relaunched from scratch. \\
11 & after the results & End of the simulations; descriptive complement to E7 (ranking by IQM in both physics, no test). \\
12 & before any new run & Consolidation: seeds 10 to 19 (160 runs), E7 at several strengths, reading at equal budget (200,000 steps). \\
13 & after the results & Intervals made reproducible: rliable ignores \texttt{random\_state}, NumPy's global generator is seeded; verdicts unchanged. \\
14 & after the results & Corrections after a critical review: date of \file{dose.py}, counts in deviations 7 and 8, SB3 2.9.2, actual steps, undeclared relaunches, seeds read, ledger. \\
15 & no effect on the results & Up to eight cores granted to the project depending on the machine's load. \\
16 & no effect on the results & Stop by the watchdog, surviving worker fixed, battery rule. \\
17 & no effect on the results & Battery condition lifted. \\
18 & after the results, before this computation & E2: the 24 probability-of-improvement intervals made simultaneous by Bonferroni (99.79\%). \\
19 & no effect on the results & Cap of two cores for the whole project; watchdog restarted. \\
20 & after the runs, before any computation on seeds 10 to 19 & Computation of the additions of deviation 12; defect of \file{e5.py} fixed (right order over all the chosen seeds). \\
21 & after the results, before this computation & E2: strict count, a pair counts only if each beats the other, settled, on some task (defect fixed). \\
22 & after the results, before this computation & E1 redone on the sixteen (program, task) pairs, seeds 0 to 19. \\
23 & after the results, before this computation & E3 with Holm's correction over four tasks, gap in points and power. \\
24 & after the results, before this computation & E7: the Hopper reversal under modified physics, tested pair by pair. \\
25 & after the results, before this computation & E5 checked out of sample: seeds 0 to 9 judged on seeds 10 to 19, and the reverse. \\
26 & after the results, before this computation & The `settled' rule calibrated on true null hypotheses (two groups of the same program). \\
27 & after the results, before this computation & B1: the factor between the script's $t$ and the seeds' $t$ is $\sqrt{n/5}/\rho$ (defect fixed). \\
28 & a record, no statistical computation & The hyperparameters actually used, read from the built programs and the saved models. \\
29 & after the results, before this computation & Learning curves: ranking during training, final score against best point, autocorrelation. \\
30 & before this computation & Effect sizes in score points (Hodges-Lehmann) for E2, E3 and E6. \\
31 & after the results, before this computation & E6: speed of the official agents read from their training logs. \\
32 & before any training run & Determinism of SAC and CrossQ: three training runs redone identically. \\
33 & before these evaluations and this training & E6: official agents replayed (speed, videos), the agent published by the Zoo evaluated, recipe with $\gamma = 0.99$. \\
34 & before any training run & Why the two PPOs diverge: seeds 20 to 39 and orthogonal initialisation swapped. \\
\addlinespace
\multicolumn{3}{@{}l}{\textbf{GRPO} (\file{sim/llm/LLM\_PROTOCOL-en.md}, 20 deviations)} \\
1 & after reading outputs, before any training & Length limit raised from 64 to 160 tokens. \\
2 & before any training & Measures reduced to 150 problems for greedy decoding and 50 for sampling, for cost. \\
3 & after the pilot, before any training & Task \texttt{mul2} chosen; main measure: sampled pass@1. \\
4 & report, no retained result & The learning-rate pilot saturates memory (swap); killed, no usable result. \\
5 & report, no retained result & Diagnosis without training: 23 answers out of 32 cut off at 160 tokens at step 10 of the pilot. \\
6 & before any retained run & Reduced version: 100 steps, seeds 0 to 2, learning rate 3e-6, updates in batches of 8 answers, memory cap. \\
7 & rule fixed before any result for the seed & Pause at our request; \texttt{nan} crash of a Qwen run; a single relaunch planned. \\
8 & report, no retained result & Resumption: none of the nine remaining runs completes (\texttt{nan}, macOS cut-off, MPS memory). \\
9 & before any result of these runs & Relaunch on an idle machine (\file{night.sh}) with a control run first. \\
10 & report, no retained result & Cause of the \texttt{nan} found: non-finite gradients on MPS, clean on CPU; recomputation safeguards. \\
11 & before any retained run with this code & The remaining Qwen runs move to CPU; Llama stays on GPU, on an idle machine. \\
12 & corrections before any relaunch & Qwen finished on CPU; the six Llama runs run out of MPS memory; cache emptied at every step, cap raised. \\
13 & statistic details set after the results & All twelve runs finished; percentile bootstrap over the seeds, 10,000 draws. \\
14 & before any new run & Five seeds: seeds 3 and 4 added for both models and both rewards. \\
15 & after the results & Corrections after a critical review (counts of deviation 8, dates); no result changes. \\
16 & after the step-100 values, before any gain & Llama, random reward, seed 4 fails twice: counted as a result; statistic over the finished runs, sensitivity to the missing gain. \\
17 & no effect on the results & The Llama CPU run starts after the queue, under the two-core cap. \\
18 & after the results & Defects found: repetition penalty of 1.1 for Qwen, pilot read on the test set, device and seed mixed. \\
19 & before the script & Analyses from the logs: greedy accuracy, pass@8, Student intervals, true reward against coin flip, truncation. \\
20 & before any run of the rerun & Rerun planned: 10 seeds, no repetition penalty, pilot off the test set; nothing is started. \\
\end{longtable}
\endgroup
